\chapter{Week 4: Names and Functions}
\label{ch:week4}

Lecture 4: September 29, 2026. Recitation 4: October 2, 2026.

Supplementary reading:
\begin{itemize}
  \item \cite{PLAI} Chapter 5 (Adding Functions to the Language), which implements functions by substitution as this week does, and Chapter 6 (From Substitution to Environments) for where the course goes next.
  \item \cite{TAPL} Sections 5.2--5.3 (the untyped lambda calculus: free variables, substitution, and $\alpha$-conversion), optional.
\end{itemize}

Lecture~4 gives the calculator \emph{names} and then \emph{functions}. Both are implemented by substitution, and the lecture's real subject is how subtle substitution is: when to evaluate the thing being substituted, how inner bindings shadow outer ones, and how a careless substitution can \emph{capture} a variable and change a program's meaning. Recitation~4 adds a \code{match} form to the same interpreter. This chapter covers:
\begin{itemize}
  \item \Cref{sec:letsubst}: \code{idC} and \code{letC}, eager substitution, and shadowing;
  \item \Cref{sec:lazy}: lazy substitution, its costs and benefits, and the variable-capture bug;
  \item \Cref{sec:capavoid}: capture-avoiding substitution with fresh names from \code{gensym};
  \item \Cref{sec:lambda}: \code{lambda} and application, a new \tyname{Value} type, and capture again;
  \item \Cref{sec:rec4}: Recitation~4, pattern matching in the interpreter.
\end{itemize}

\section{Names: \texttt{let} and Substitution}
\label{sec:letsubst}

\subsection{The Plan}

At the end of \Cref{ch:week3} the calculator had \code{numC}, \code{plusC}, and \code{timesC}. The first new form is \code{(let x e1 e2)}: compute \code{e1}, replace every \code{x} in \code{e2} with that value, then compute \code{e2}.
\begin{lstlisting}[style=repl]
(let x 5 (+ x x))
---> (+ 5 5)                 replace every x in the body with 5
---> 10

(let x (+ 1 5) (+ x x))
---> (let x 6 (+ x x))       e1 is evaluated first
---> (+ 6 6)
---> 12
\end{lstlisting}

\subsection{Two New Variants}

\begin{lstlisting}
(define-type ArithE
  [numC (n : number)]
  [idC (x : symbol)]
  [plusC (e1 : ArithE) (e2 : ArithE)]
  [timesC (e1 : ArithE) (e2 : ArithE)]
  [letC (x : symbol) (e1 : ArithE) (e2 : ArithE)])
\end{lstlisting}
(\Cref{sec:lambda} adds two more variants for functions; from then on every function on \tyname{ArithE} has a case for them.)

\subsubsection{What is \code{idC} for?} An \code{idC} is a \emph{use} of a name, such as the \code{x} in \code{(+ x 1)}. The operands of \code{plusC} must be \tyname{ArithE}s, so a bare symbol \code{'x} cannot go there; \code{idC} wraps the symbol so a name can sit in the tree until substitution replaces it. It does for symbols what \code{numC} does for numbers.

A name appears in two roles, and only one of them is an \code{idC}:
\begin{lstlisting}
;; (let x 5 (+ x 1))
(letC 'x (numC 5) (plusC (idC 'x) (numC 1)))
;;    ^ binding: a plain symbol    ^ use: an idC
\end{lstlisting}
The name right after \code{let} is being \emph{bound}. It is a plain symbol field of \code{letC}, because it is never evaluated or replaced; it only tells substitution which name to look for. Symbols like \code{+} and \code{let} never become \code{idC}s either: \code{parse} consumes them to decide which variant to build. Writing \rkt|(plusC 'x (numC 1))| is a type error, since \code{'x} is a \tyname{symbol} and not an \tyname{ArithE}.

\subsubsection{Parsing} A lone symbol becomes an \code{idC}. In \code{(let x e1 e2)} the \code{x} is taken as a plain symbol, while \code{e1} and \code{e2} are parsed recursively:
\begin{lstlisting}
(define (parse (s : s-expression)) : ArithE
  (cond
    [(s-exp-number? s) (numC (s-exp->number s))]
    [(s-exp-symbol? s) (idC (s-exp->symbol s))]
    [(s-exp-list? s)
     (let [(l (s-exp->list s))]
       (case (s-exp->symbol (first l))
         [(+) (plusC (parse (second l)) (parse (third l)))]
         [(*) (timesC (parse (second l)) (parse (third l)))]
         [(let) (letC (s-exp->symbol (second l)) (parse (third l)) (parse (fourth l)))]))]))
\end{lstlisting}
As in lecture~3, there is no error checking: \rkt|(parse '(let x 1))| fails inside \code{fourth} with \code{list contains too few elements}.

\subsection{Eager Substitution}

\rkt|(subst-eager v x e)| replaces every \code{x} in \code{e} with the number \code{v}. Most variants just recurse into their children; the interesting cases are the ones that mention names.
\begin{lstlisting}
(define (subst-eager (v : number) (x : symbol) (e : ArithE)) : ArithE
  (type-case ArithE e
    [numC (n) (numC n)]
    [idC (y) (if (equal? x y) (numC v) (idC y))]
    [plusC (e1 e2) (plusC (subst-eager v x e1) (subst-eager v x e2))]
    [timesC (e1 e2) (timesC (subst-eager v x e1) (subst-eager v x e2))]
    [letC (y e1 e2)
          (letC y
                (subst-eager v x e1)
                (if (equal? x y) e2 (subst-eager v x e2)))]))
\end{lstlisting}

\subsubsection{Why \code{(numC v)} and \code{(idC y)}?} \code{subst-eager} returns an \tyname{ArithE} tree, but \code{v} is a number, so it is wrapped as \rkt|(numC v)| to go back into the tree. An \code{idC} for a different name is left alone: \rkt|(idC y)| rebuilds the same node, which is equivalent to returning \code{e}.
\begin{lstlisting}
(subst-eager 5 'x (plusC (idC 'x) (idC 'y)))   ; = (plusC (numC 5) (idC 'y))
\end{lstlisting}

\subsection{Shadowing}

What is \code{(let x 5 (let x 6 (+ x x)))}: 12 or 10? The answer is 12. The inner \code{let} binds its own \code{x}, so inside its body \code{x} means 6: the inner \code{x} \emph{shadows} the outer one (\Cref{fig:shadow}).

\begin{figure}[htbp]
\centering
\begin{tikzpicture}[tok/.style={inner sep=0pt, font=\ttfamily\normalsize, anchor=base west}]
  \node[tok] at (0,0) {(let x 5 (let x 6 (+ x x)))};
  \node[tok] (x1) at (5\cw,0) {\phantom{x}};
  \node[tok] (x2) at (14\cw,0) {\phantom{x}};
  \node[tok] (u1) at (21\cw,0) {\phantom{x}};
  \node[tok] (u2) at (23\cw,0) {\phantom{x}};
  \draw[bind] (u1.south) to[out=-110, in=-70] (x2.south);
  \draw[bind] (u2.south) to[out=-110, in=-70] (x2.south);
  \node[lbl, below=8mm of u1, text=accent] {both uses refer to the inner binder};
  \node[lbl, above=1mm of x1] {outer};
  \node[lbl, above=1mm of x2] {inner};
\end{tikzpicture}
\caption{Each use of a name refers to the nearest enclosing binder of that name.}
\label{fig:shadow}
\end{figure}

That is the \rkt|(equal? x y)| check in the \code{letC} case: when an inner \code{let} binds the same name we are substituting for, substitution stops at its body.
\begin{lstlisting}[style=repl]
(let x 5 (let x 6 (+ x x)))
---> (let x 6 (+ x x))       the outer subst leaves the inner body alone
---> (+ 6 6)
---> 12
\end{lstlisting}
Substituting blindly would give \code{(let x 6 (+ 5 5))}, which is 10. The inner \code{e1} is still substituted, because it sits \emph{outside} the inner binding; the inner \code{x} exists only in the inner body:
\begin{lstlisting}[style=repl]
(let x 1 (let x (+ x 1) x))
---> (let x (+ 1 1) x)       the x in e1 is the outer x
---> 2
\end{lstlisting}

\subsection{\texttt{eval-eager}}

\begin{lstlisting}
(define (eval-eager (e : ArithE)) : number
  (type-case ArithE e
    [numC (n) n]
    [idC (x) (error 'eval "unbound variable")]
    [plusC (e1 e2) (+ (eval-eager e1) (eval-eager e2))]
    [timesC (e1 e2) (* (eval-eager e1) (eval-eager e2))]
    [letC (y e1 e2) (eval-eager (subst-eager (eval-eager e1) y e2))]))
\end{lstlisting}
The \code{letC} line is the plan written out: evaluate \code{e1}, substitute the resulting number for \code{y} in \code{e2}, and evaluate the result. \code{eval-eager} never meets an \code{idC} that some \code{let} bound, because the enclosing \code{letC} has already substituted it away. An \code{idC} that survives was never bound, as in \code{(+ x 1)}, so it is an error.
\begin{lstlisting}[style=repl]
(let x 5 (let y (+ x 1) (* x y)))
letC x: eval e1 = 5, subst 5 for x in the body
---> (let y (+ 5 1) (* 5 y))
letC y: eval e1 = 6, subst 6 for y in the body
---> (* 5 6)
---> 30
\end{lstlisting}
Each \code{let} removes the \code{idC}s for its own name; once every \code{let} has been processed, none is left.

\section{Lazy Substitution and Variable Capture}
\label{sec:lazy}

\subsection{Lazy Substitution}

An alternative: do not evaluate \code{e1} first. Substitute the unevaluated tree into the body, where it is evaluated each time the body uses it.
\begin{lstlisting}[style=repl]
(let x (+ 1 1) (+ x 2))
eager:  ---> (let x 2 (+ x 2))  ---> (+ 2 2)  ---> 4
lazy:   ---> (+ (+ 1 1) 2)                    ---> 4
\end{lstlisting}
\code{subst-lazy} is \code{subst-eager} with \code{v} an \tyname{ArithE}, the unevaluated \code{e1}, which goes into the tree as is, with no \code{numC} wrapping:
\begin{lstlisting}
(define (subst-lazy (v : ArithE) (x : symbol) (e : ArithE)) : ArithE
  (type-case ArithE e
    [numC (n) e]
    [idC (y) (if (equal? x y) v e)]
    [plusC (e1 e2) (plusC (subst-lazy v x e1) (subst-lazy v x e2))]
    [timesC (e1 e2) (timesC (subst-lazy v x e1) (subst-lazy v x e2))]
    [letC (y e1 e2)
          (letC y
                (subst-lazy v x e1)
                (if (equal? x y) e2 (subst-lazy v x e2)))]))

(define (eval-lazy (e : ArithE)) : number
  (type-case ArithE e
    [numC (n) n]
    [idC (x) (error 'eval "unbound variable")]
    [plusC (e1 e2) (+ (eval-lazy e1) (eval-lazy e2))]
    [timesC (e1 e2) (* (eval-lazy e1) (eval-lazy e2))]
    [letC (y e1 e2) (eval-lazy (subst-lazy e1 y e2))]))
\end{lstlisting}

\begin{pitfall}[Mixing evaluators]
Every recursive call in \code{eval-lazy} must be to \code{eval-lazy}. If one called \code{eval-eager}, the type checker would not complain, since both return a number, but inner \code{let}s would silently become eager.
\end{pitfall}

\subsection{Is Lazy More Efficient?}

When \code{x} is unused, yes: \code{e1} is never evaluated.
\begin{lstlisting}[style=repl]
(let x (+ 5 6) 10)
eager:  (let x 11 10) ---> 10      (+ 5 6) computed once, for nothing
lazy:   10                         (+ 5 6) never computed
\end{lstlisting}
So \code{(let x (/ 5 0) 10)} would not be an error in the lazy version (if the language had \code{/}), and neither is an unbound name in \code{e1}: \code{(let x (+ y y) 10)} is an error eagerly but 10 lazily.

When \code{x} is used more than once, lazy is \emph{less} efficient, since \code{e1} is recomputed at every use:
\begin{lstlisting}[style=repl]
(let x (+ 5 6) (+ (+ x x) x))
eager:  (+ 5 6) once     ---> (+ (+ 11 11) 11)
lazy:   (+ 5 6) 3 times  ---> (+ (+ (+ 5 6) (+ 5 6)) (+ 5 6))
\end{lstlisting}

\subsection{Where Eager and Lazy Differ}

\begin{center}
\begin{tabular}{@{}lll@{}}
\toprule
program & eager & lazy \\
\midrule
\code{(let x (+ y y) 10)} & \code{unbound variable} & \code{10} \\
\code{(let y x (let x 2 y))} & \code{unbound variable} & \code{2} \\
\bottomrule
\end{tabular}
\end{center}
The first difference is harmless: \code{x} is unused, so lazy never evaluates \code{e1} and never notices the unbound \code{y}. The second is a real bug.

\subsection{Variable Capture}

\begin{lstlisting}[style=repl]
(let x y (let y 2 (+ x y)))
lazy:
---> (let y 2 (+ y y))        the free y is captured by the inner let
---> (+ 2 2)
---> 4
\end{lstlisting}
The outer \code{y} is free: nothing binds it, so the correct answer is an error. Racket agrees; the same program in \code{\#lang racket} is rejected with \code{y: unbound identifier} before it runs. But lazy substitution moved that free \code{y} underneath \code{(let y 2 ...)}, where it now refers to the inner \code{y} (\Cref{fig:capture}).

\begin{figure}[htbp]
\centering
\begin{tikzpicture}[tok/.style={inner sep=0pt, font=\ttfamily\normalsize, anchor=base west}]
  \node[tok] at (0,0) {(let x y (let y 2 (+ x y)))};
  \node[tok] at (0,-2.2) {(let y 2 (+ y y))};
  \node[tok] (fy) at (7\cw,0) {\phantom{x}};
  \node[tok] (by) at (14\cw,0) {\phantom{x}};
  \node[tok] (uy) at (23\cw,0) {\phantom{x}};
  \draw[bind] (uy.south) to[out=-110, in=-70] (by.south);
  \node[lbl, above=1mm of fy, text=sigc] {free};
  \node[tok] (by2) at (5\cw,-2.2) {\phantom{x}};
  \node[tok] (cy) at (12\cw,-2.2) {\phantom{x}};
  \node[tok] (uy2) at (14\cw,-2.2) {\phantom{x}};
  \draw[bind] (uy2.south) to[out=-110, in=-70] (by2.south);
  \draw[bad] (cy.south) to[out=-120, in=-60] (by2.south);
  \node[lbl, above=1mm of cy, text=sigc] {was free, now captured};
  \draw[arr] (10\cw,-0.75) -- (10\cw,-1.6) node[midway, right, lbl] {\code{subst-lazy (idC 'y) 'x} on the body};
\end{tikzpicture}
\caption{Capture. The substituted \code{y} (red arrow) ends up bound by an inner \code{let} it had nothing to do with.}
\label{fig:capture}
\end{figure}

Now rename the inner \code{y} to \code{z}:
\begin{lstlisting}[style=repl]
(let x y (let z 2 (+ x z)))
lazy:
---> (let z 2 (+ y z))
---> (+ y 2)
---> eval: unbound variable
\end{lstlisting}

\begin{keypoint}[Names of bound variables should not matter]
The two programs differ only in renaming the inner \code{y} to \code{z}, so they must mean the same thing. Yet lazy substitution gives 4 for one and an error for the other. This bug is called \emph{variable capture}, and a substitution that can cause it is \emph{capturing}.
\end{keypoint}

\code{eval-eager} does not have this problem: it only ever substitutes numbers, and a number contains no names to capture. Lazy substitution substitutes whole expressions, which can contain free variables. (Once values can be functions, eager evaluation substitutes code too, and capture returns; see \Cref{sec:eagerapp}.)

\subsection{(SUPPLEMENT) Further Topics}

\subsubsection{Evaluation strategies} The two substitution regimes are the two classic \emph{parameter-passing strategies}. Eager substitution is \emph{call-by-value}, used by Racket, OCaml, Python, Java, and most languages. Lazy substitution is \emph{call-by-name}, used for procedure parameters in Algol~60. The recomputation problem has a standard fix, \emph{call-by-need}: substitute a reference to a shared \emph{thunk} that evaluates \code{e1} the first time it is used and remembers the result. That is Haskell's evaluation strategy, and it keeps lazy's advantage (unused arguments are never evaluated) without its cost. Call-by-need gives the same results as call-by-name in a pure language; with effects, the number of times an effect happens becomes observable.

\subsubsection{Laziness is visible} A strategy is not just an efficiency choice when evaluation can fail or loop. If \code{e1} loops forever, \code{(let x e1 10)} loops under call-by-value and returns 10 under call-by-name. This lets lazy languages define their own control structures as ordinary functions: \code{if} could be a function in Haskell, but not in Racket, where both branches would be evaluated before the call.

\subsubsection{Capture in practice} Capture is not only a toy-interpreter bug. C preprocessor macros substitute text and capture freely: a macro body that declares \code{int tmp} silently changes the meaning of a caller's \code{tmp}. Racket's \emph{hygienic} macros rename binders automatically to prevent exactly this, which is the macro-level version of the fix in the next section.

\section{Capture-Avoiding Substitution}
\label{sec:capavoid}

\subsection{Rename, Then Substitute}

The fix is \emph{capture-avoiding substitution}: before substituting into the body of an inner \code{let}, rename its bound variable to a fresh name that appears nowhere else. Renaming a bound variable never changes what a program means (that was the point of the \code{y}/\code{z} example), so we are free to do it whenever it helps.
\begin{lstlisting}[style=repl]
(let x y (let y 2 (+ x y)))
---> (let x y (let g3189 2 (+ x g3189)))    rename the inner y first
---> (let g3189 2 (+ y g3189))              the free y is no longer captured
---> (+ y 2)
---> eval: unbound variable                 same answer as the z version
\end{lstlisting}
The only question is where the fresh name comes from.

\subsection{\texttt{gensym}: Fresh Names}

\code{gensym} (``generate symbol'') is a Racket function that returns a brand-new symbol every time it is called:
\begin{lstlisting}[style=repl]
> (gensym)
'g3189
> (gensym)
'g3190
> (equal? (gensym) (gensym))
#f
\end{lstlisting}
A name chosen by hand, such as \code{'tmp}, does not work: the program might already use \code{tmp}, and then the renaming would create exactly the capture it was meant to avoid.

\begin{definition}[Uninterned symbol]
An ordinary symbol is \emph{interned}: \code{'y} is the same symbol wherever it is written, so \rkt|(equal? 'y 'y)| is \code{\#t}. The symbol returned by \code{gensym} is \emph{uninterned}: it is different from every other symbol, including one spelled the same way.
\end{definition}
Even if a program literally contained the name \code{g3189}, it would not be \code{equal?} to the \code{g3189} that \code{gensym} returned:
\begin{lstlisting}
(define g (gensym))                                ; prints as 'g3189
(equal? g (string->symbol (symbol->string g)))     ; #f
\end{lstlisting}
So a gensym can never be captured by, or capture, any name the user wrote.

\subsubsection{Importing it into \code{plai-typed}} \code{plai-typed} does not provide \code{gensym}, and it cannot check a plain Racket import without knowing its type. \code{typed-in} imports a function from a Racket module and declares its type at the same time:
\begin{lstlisting}
(require (typed-in racket/base [gensym : (-> symbol)]))
\end{lstlisting}
\code{(-> symbol)} means ``takes no arguments and returns a symbol''. The checker trusts this declaration; it does not check Racket's code.

\subsection{\texttt{subst-lazy-v2}}

Only the binding case changes. If the inner \code{let} binds the same name as the one being replaced, stop, as before. Otherwise:
\begin{enumerate}
  \item let \code{z} be \rkt|(gensym)|, a name that appears nowhere else;
  \item rename \code{y} to \code{z} in the body: \rkt|(subst-lazy (idC z) y e2)|. This cannot cause capture, because no \code{let} inside \code{e2} binds the fresh \code{z};
  \item substitute \code{v} for \code{x} in the renamed body.
\end{enumerate}
\begin{lstlisting}
(define (subst-lazy-v2 (v : ArithE) (x : symbol) (e : ArithE)) : ArithE
  (type-case ArithE e
    [numC (n) e]
    [idC (y) (if (equal? x y) v e)]
    [plusC (e1 e2) (plusC (subst-lazy-v2 v x e1) (subst-lazy-v2 v x e2))]
    [timesC (e1 e2) (timesC (subst-lazy-v2 v x e1) (subst-lazy-v2 v x e2))]
    [letC (y e1 e2)
          (if (equal? x y)
              (letC y (subst-lazy-v2 v x e1) e2)
              (let [(z (gensym))]
                (letC z
                      (subst-lazy-v2 v x e1)
                      (subst-lazy-v2 v x (subst-lazy (idC z) y e2)))))]))
\end{lstlisting}
\begin{lstlisting}
(subst-lazy-v2 (idC 'y) 'x (parse '(let y 2 (+ x y))))
;; = (letC 'g3445 (numC 2) (plusC (idC 'y) (idC 'g3445)))
\end{lstlisting}
The inner \code{y} became \code{g3445}, and the \code{y} that came from outside is still free, so evaluation reports it as unbound. Now the \code{y} and \code{z} versions of the program agree.

The lecture's version (\code{lec4a.rkt}) splits the \code{letC} case three ways: \code{x} equals \code{y}; \code{y} is \code{possibly-free?} in \code{v}, so rename; otherwise substitute without renaming. Its \code{possibly-free?} is a stub that always returns \code{\#t}, so it behaves exactly like the two-way code above. Implementing \code{possibly-free?} properly, so that renaming happens only when needed, is left as an exercise (see the free-variable definition in \Cref{sec:letsubst}).

\begin{remark}
Step~2 uses the old \code{subst-lazy}, which may capture in general. It is safe here because the expression being substituted, \rkt|(idC z)|, contains only the fresh name \code{z}, which nothing in \code{e2} can bind.
\end{remark}

\section{Functions: \texttt{lambda} and Application}
\label{sec:lambda}

\subsection{Syntax}

The last part of the lecture adds functions:
\begin{center}
\begin{tabular}{@{}lll@{}}
\toprule
syntax & meaning & abstract syntax \\
\midrule
\code{(lambda x e)} & a function with parameter \code{x} and body \code{e} & \code{lambdaC} \\
\code{(e1 e2)} & call the function \code{e1} on the argument \code{e2} & \code{appC} \\
\bottomrule
\end{tabular}
\end{center}
\begin{lstlisting}
(define-type ArithE
  [numC (n : number)]
  [idC (x : symbol)]
  [plusC (e1 : ArithE) (e2 : ArithE)]
  [timesC (e1 : ArithE) (e2 : ArithE)]
  [letC (x : symbol) (e1 : ArithE) (e2 : ArithE)]
  [lambdaC (x : symbol) (e : ArithE)]
  [appC (e1 : ArithE) (e2 : ArithE)])
\end{lstlisting}
\code{lambdaC} \emph{defines} a function and \code{appC} \emph{uses} one, much as \code{letC} introduces a name and \code{idC} uses it. Functions take one argument; several arguments can be simulated by currying (\Cref{sec:hof}).

\subsubsection{Parsing} Any list whose head is not \code{+}, \code{*}, \code{let}, or \code{lambda} is an application:
\begin{lstlisting}
(define (parse (s : s-expression)) : ArithE
  (cond
    [(s-exp-number? s) (numC (s-exp->number s))]
    [(s-exp-symbol? s) (idC (s-exp->symbol s))]
    [(s-exp-list? s)
     (let [(l (s-exp->list s))]
       (if (s-exp-symbol? (first l))
           (case (s-exp->symbol (first l))
             [(+) (plusC (parse (second l)) (parse (third l)))]
             [(*) (timesC (parse (second l)) (parse (third l)))]
             [(let) (letC (s-exp->symbol (second l)) (parse (third l)) (parse (fourth l)))]
             [(lambda) (lambdaC (s-exp->symbol (second l)) (parse (third l)))]
             [else (appC (parse (first l)) (parse (second l)))])
           ;; head is not a symbol, e.g. ((lambda x (+ x 1)) 5)
           (appC (parse (first l)) (parse (second l)))))]))
\end{lstlisting}
\begin{lstlisting}
(parse '((lambda x (+ x 1)) 5))
;; = (appC (lambdaC 'x (plusC (idC 'x) (numC 1))) (numC 5))
(parse '(f 5))
;; = (appC (idC 'f) (numC 5))
\end{lstlisting}
The outer \code{if} matters for calls whose head is itself a list, like \code{((lambda x (+ x 1)) 5)}: calling \code{s-exp->symbol} on it would fail with \code{s-exp->symbol: not a symbol}.

\begin{remark}
With an application fallback, \code{'(- 1 2)} no longer fails to parse. The \code{else} branch reads it as a call of a function named \code{-}, giving \rkt|(appC (idC '-) (numC 1))| and silently dropping the \code{2}. Before \code{lambda} was added, it was \code{case: no matching clause}.
\end{remark}

\subsection{\texttt{lambda} in Substitution}

A \code{lambda} binds its parameter in its body, exactly as \code{let} binds \code{x} in \code{e2}. So the \code{lambdaC} case of each substitution function is its \code{letC} case without an \code{e1}: stop at shadowing, and in \code{subst-lazy-v2} rename with \code{gensym} first. The application case just recurses.
\begin{lstlisting}
;; in subst-eager (subst-lazy is the same with subst-lazy)
[lambdaC (y e)
         (if (equal? x y) (lambdaC y e) (lambdaC y (subst-eager v x e)))]
[appC (e1 e2) (appC (subst-eager v x e1) (subst-eager v x e2))]

;; in subst-lazy-v2
[lambdaC (y e)
         (if (equal? x y)
             (lambdaC y e)
             (let [(z (gensym))]
               (lambdaC z (subst-lazy-v2 v x (subst-lazy (idC z) y e)))))]
[appC (e1 e2) (appC (subst-lazy-v2 v x e1) (subst-lazy-v2 v x e2))]
\end{lstlisting}

\begin{pitfall}[The pattern variable shadows the argument]
In \rkt|[lambdaC (y e) ...]|, the pattern variable \code{e} (the body) shadows \code{subst}'s own argument \code{e} (the whole \code{lambda}). So the shadowing branch must rebuild \rkt|(lambdaC y e)|; returning plain \code{e} would drop the \code{lambda} and return only its body. Shadowing bites the interpreter's author too.
\end{pitfall}

\subsection{A New Result Type}

A program can now produce a function, for example \code{(lambda x (+ x 1))} or \code{(let f (lambda x x) f)}, and a function is not a number. So evaluation returns a \tyname{Value}, which is either:
\begin{lstlisting}
(define-type Value
  [numV (n : number)]
  [funV (x : symbol) (e : ArithE)])
\end{lstlisting}
\code{+} and \code{*} only make sense on numbers, so they unwrap \code{numV} first and complain about anything else:
\begin{lstlisting}
(define (num+ (v1 : Value) (v2 : Value)) : Value
  (if (and (numV? v1) (numV? v2))
      (numV (+ (numV-n v1) (numV-n v2)))
      (error 'eval "not a number")))

(define (num* (v1 : Value) (v2 : Value)) : Value
  (if (and (numV? v1) (numV? v2))
      (numV (* (numV-n v1) (numV-n v2)))
      (error 'eval "not a number")))
\end{lstlisting}
This is the first \emph{dynamic} type error in our language: whether \code{(+ e1 e2)} is meaningful depends on what \code{e1} and \code{e2} evaluate to, which the interpreter only learns at run time.

\subsection{Eager Application}
\label{sec:eagerapp}

The lecture first adds functions to the \emph{eager} interpreter (\code{lec4b.rkt}, \code{lec4c.rkt}). Application is \emph{call-by-value}: evaluate the function, evaluate the argument, then substitute the argument's \emph{value} into the body. \code{let} works the same way.

\code{subst-eager} substitutes only numbers, but a value may now be a function. So the value is turned back into an expression, and substituted with the expression-substituting \code{subst-lazy} or \code{subst-lazy-v2}. The expression they receive is already evaluated, so nothing is lazy about it:
\begin{lstlisting}
(define (value->expr (v : Value)) : ArithE
  (type-case Value v
    [numV (n) (numC n)]
    [funV (x e) (lambdaC x e)]))

(define (eval-eager-v2 (e : ArithE)) : Value
  (type-case ArithE e
    [numC (n) (numV n)]
    [idC (x) (error 'eval "unbound variable")]
    [plusC (e1 e2) (num+ (eval-eager-v2 e1) (eval-eager-v2 e2))]
    [timesC (e1 e2) (num* (eval-eager-v2 e1) (eval-eager-v2 e2))]
    [letC (y e1 e2)
          (eval-eager-v2 (subst-lazy (value->expr (eval-eager-v2 e1)) y e2))]
    [lambdaC (x e) (funV x e)]
    [appC (e1 e2)
          (type-case Value (eval-eager-v2 e1)
            [funV (x e)
                  (eval-eager-v2 (subst-lazy (value->expr (eval-eager-v2 e2)) x e))]
            [numV (n) (error 'eval "not a function")])]))
\end{lstlisting}
The lecture's \code{lec4b.rkt} writes this with \code{Expr} for \tyname{ArithE} and \code{lambdaV} for \code{funV}, and gives \code{subst} a \tyname{Value} that it converts at the \code{idC} case; the effect is the same. Functions can now be passed around:
\begin{lstlisting}[style=repl]
(let do-twice (lambda f (f (f 1)))
  (let double (lambda x (* 2 x))
    (do-twice double)))
---> (do-twice double)  with both substituted
---> (double (double 1))
---> 4
\end{lstlisting}

\subsubsection{Eager capture} \Cref{sec:lazy} said that eager substitution cannot capture because a number has no names in it. A function value does have names, and its body may contain free variables:
\begin{lstlisting}[style=repl]
(let f (lambda x (+ x y))         y is free in f's body
  (let y 6
    (f 1)))
---> (let y 6 ((lambda x (+ x y)) 1))   substitute f's value: y is captured
---> ((lambda x (+ x 6)) 1)
---> 7
\end{lstlisting}
The right answer is an error: the \code{y} in \code{f} refers to no binding, and Racket rejects the same program with \code{y: unbound identifier}. \code{eval-eager-v2} returns 7 because \code{subst-lazy} carried the free \code{y} under \code{(let y 6 ...)}. Eager evaluation is not a defense against capture once values can contain code.

The fix (\code{lec4c.rkt}) is the one from \Cref{sec:capavoid}: \code{eval-eager-v3} is \code{eval-eager-v2} with both calls to \code{subst-lazy} replaced by \code{subst-lazy-v2}. Now the program above raises \code{unbound variable}, and returning \code{f} instead of calling it shows the renaming:
\begin{lstlisting}[style=repl]
> (eval-eager-v3 (parse '(let f (lambda x (+ x y)) (let y 6 f))))
(funV 'g3507 (plusC (idC 'g3507) (idC 'y)))
\end{lstlisting}
Substituting \code{f} into the body of \code{(let y 6 ...)} renamed \code{f}'s parameter to a fresh name, as the \code{lambdaC} case of \code{subst-lazy-v2} always does; the free \code{y} stays free.

\subsection{Lazy Application: \texttt{eval-lazy-v2}}

The last version (\code{lec4d.rkt}) makes application lazy, like \code{let} in \Cref{sec:lazy}: the argument is substituted unevaluated.

\begin{lstlisting}
(define (eval-lazy-v2 (e : ArithE)) : Value
  (type-case ArithE e
    [numC (n) (numV n)]
    [idC (x) (error 'eval "unbound variable")]
    [plusC (e1 e2) (num+ (eval-lazy-v2 e1) (eval-lazy-v2 e2))]
    [timesC (e1 e2) (num* (eval-lazy-v2 e1) (eval-lazy-v2 e2))]
    [letC (y e1 e2) (eval-lazy-v2 (subst-lazy-v2 e1 y e2))]
    [lambdaC (x e) (funV x e)]
    [appC (e1 e2)
          (type-case Value (eval-lazy-v2 e1)
            [funV (x e) (eval-lazy-v2 (subst-lazy-v2 e2 x e))]
            [numV (n) (error 'eval "not a function")])]))
\end{lstlisting}
\begin{itemize}
  \item \textbf{\code{lambdaC}}: a \code{lambda} is already a value. There is nothing to compute until it is called.
  \item \textbf{\code{appC}}: evaluate \code{e1}; it must be a function \rkt|(funV x e)|. Then call it the way \code{let} works: substitute the argument \code{e2} for \code{x} in the body \code{e}, and evaluate the result. \code{e2} is substituted unevaluated, because this is the lazy version.
\end{itemize}
\begin{lstlisting}[style=repl]
((lambda x (+ x 1)) 5)
---> (+ 5 1)             subst 5 for x in the body
---> 6
\end{lstlisting}

\begin{keypoint}[\code{let} is application]
\code{((lambda x e) e2)} means the same as \code{(let x e2 e)}. What \code{lambda} adds is that the function is a \emph{value}: it can be named with \code{let}, passed as an argument, and returned as a result.
\end{keypoint}

\begin{center}
\begin{tabular}{@{}ll@{}}
\toprule
program & result \\
\midrule
\code{(let f (lambda x (* x 2)) (f 5))} & \code{(numV 10)} \\
\code{((lambda f (f 3)) (lambda x (+ x 1)))} & \code{(numV 4)}: a function as argument \\
\code{(((lambda x (lambda y (+ x y))) 1) 2)} & \code{(numV 3)}: a function as result \\
\code{((lambda x 5) y)} & \code{(numV 5)}: the unused argument is never evaluated \\
the \code{do-twice} example of \Cref{sec:eagerapp} & \code{(numV 4)} \\
\code{(5 3)} & \code{not a function} \\
\code{(+ (lambda x x) 1)} & \code{not a number} \\
\bottomrule
\end{tabular}
\end{center}

\subsection{Capture, Again}

With functions, capture-avoiding substitution matters even more. A function can return a function, so a free name can end up under another \code{lambda}:
\begin{lstlisting}[style=repl]
(((lambda x (lambda y (+ x y))) y) 2)
subst-lazy:     (lambda y (+ y y)),  then 2 ---> 4          captured!
subst-lazy-v2:  (lambda g (+ y g)),  then 2 ---> (+ y 2) ---> unbound variable
\end{lstlisting}
\code{eval-eager} and \code{eval-lazy} still return a \code{number}, and \code{subst-eager} can substitute only numbers, so their \code{lambdaC} and \code{appC} cases remain errors (\code{"TODO: lambda"}); \code{eval-eager-v2}/\code{v3} and \code{eval-lazy-v2} are the versions with functions.

\begin{center}
\begin{tabular}{@{}lll@{}}
\toprule
program & \code{eval-eager-v3} (call-by-value) & \code{eval-lazy-v2} (call-by-name) \\
\midrule
\code{((lambda x 5) y)} & \code{unbound variable} & \code{(numV 5)} \\
\code{(let x (+ y y) 10)} & \code{unbound variable} & \code{(numV 10)} \\
the \code{f}/\code{y} example of \Cref{sec:eagerapp} & \code{unbound variable} & \code{unbound variable} \\
\bottomrule
\end{tabular}
\end{center}
With capture-avoiding substitution, the two strategies differ only where an argument is never used, as for \code{let} in \Cref{sec:lazy}. Racket, like most languages, is call-by-value.

\section{Recitation 4: Pattern Matching in the Interpreter}
\label{sec:rec4}

In \Cref{sec:rec3} we \emph{used} Racket's \code{match} to write \code{factorial} and \code{fibonacci}. This recitation \emph{implements} a small version of it: the calculator with \code{let} gets a \code{match} form, and each part of the interpreter is extended in turn: the types, \code{parse}, \code{subst}, and \code{eval}. The file starts with the \code{gensym} import of \Cref{sec:capavoid}, since its \code{let} uses capture-avoiding lazy substitution.

\subsection{The New Form}

A pattern-matching expression looks like
\begin{lstlisting}[style=repl]
(match e (n1 e1) (n2 e2) ... (_ en))
\end{lstlisting}
where each \code{ni} is a number literal and \code{\_} is the wildcard. To evaluate it:
\begin{enumerate}
  \item evaluate \code{e} to a number, once;
  \item try the clauses top to bottom: \code{ni} matches if it equals that number, and \code{\_} matches anything;
  \item the \emph{first} matching clause wins, and only its expression is evaluated;
  \item if no clause matches, it is an error.
\end{enumerate}
This is Racket's \code{match} restricted to literals and \code{\_}, just enough for \code{factorial}- and \code{fibonacci}-style definitions. There are no pattern variables, so a clause never binds a name; that matters for substitution.

\subsection{Types}

A clause pairs a pattern with an expression. The two kinds of clause carry different data (a number pattern has a number, the wildcard nothing), so each gets its own variant:
\begin{lstlisting}
(define-type PatternE
  [numP (n : number) (e : ArithE)]
  [anyP (e : ArithE)])

(define-type ArithE
  [numC (n : number)]
  [plusC (e1 : ArithE) (e2 : ArithE)]
  [timesC (e1 : ArithE) (e2 : ArithE)]
  [letC (x : symbol) (e1 : ArithE) (e2 : ArithE)]
  [idC (x : symbol)]
  [patternC (e : ArithE) (patterns : (listof PatternE))])
\end{lstlisting}
\code{patternC} holds the expression being matched, the \emph{scrutinee}, and the clauses in order. The two types are \emph{mutually recursive}: an \tyname{ArithE} can contain \tyname{PatternE}s, and each \tyname{PatternE} contains an \tyname{ArithE}. \code{define-type} allows that, in either order.

\subsubsection{Examples} The handout's four examples, in ML-like notation:
\begin{center}
\begin{tabular}{@{}lll@{}}
\toprule
 & program & expected \\
\midrule
1 & \code{match 5 with | 5 -> 0 | 6 -> 1} & 0 \\
2 & \code{match (1 + 2) with | 3 -> 1 | \_ -> x} & 1 \\
3 & \code{match x with | 0 -> z | \_ -> (match y with | 0 -> x | 1 -> y)} & unbound variable \\
4 & \code{let x = 1 in let y = 1 in} (the \code{match} of 3, last clause \code{y + 1}) & 2 \\
\bottomrule
\end{tabular}
\end{center}
\begin{lstlisting}
(define example-2 : ArithE
  (patternC (plusC (numC 1) (numC 2))
            (list (numP 3 (numC 1))
                  (anyP (idC 'x)))))

(define example-4 : ArithE
  (letC 'x (numC 1)
        (letC 'y (numC 1)
              (patternC (idC 'x)
                        (list (numP 0 (idC 'z))
                              (anyP (patternC (idC 'y)
                                              (list (numP 0 (idC 'x))
                                                    (numP 1 (plusC (numC 1) (idC 'y)))))))))))
\end{lstlisting}
Examples 2 and 4 both contain a name that nothing binds (\code{x} and \code{z}), yet neither is an error, because that clause is never chosen. In example~3 the error comes from the scrutinee \code{x}, which is always evaluated.

\subsection{Parsing}

The concrete syntax is \code{(match e clause ...)}, so after the \code{match} keyword, \rkt|(second l)| is the scrutinee and \rkt|(rest (rest l))| is the list of clauses. A helper turns that list of s-expressions into a list of \tyname{PatternE}s. Each clause is a two-element list \code{(pattern expr)}, and the pattern decides which variant to build:
\begin{lstlisting}
(define (parse-patterns (s-pats : (listof s-expression))) : (listof PatternE)
  (cond
    [(empty? s-pats) empty]
    [else
     (let [(s-pat (s-exp->list (first s-pats)))]
       (cond
         [(s-exp-number? (first s-pat))
          (cons (numP (s-exp->number (first s-pat)) (parse (second s-pat)))
                (parse-patterns (rest s-pats)))]
         [(equal? (s-exp->symbol (first s-pat)) '_)
          (list (anyP (parse (second s-pat))))]))]))

(define (parse (s : s-expression)) : ArithE
  (cond
    [(s-exp-number? s) (numC (s-exp->number s))]
    [(s-exp-symbol? s) (idC (s-exp->symbol s))]
    [(s-exp-list? s)
     (let [(l (s-exp->list s))]
       (case (s-exp->symbol (first l))
         [(+) (plusC (parse (second l)) (parse (third l)))]
         [(*) (timesC (parse (second l)) (parse (third l)))]
         [(let) (letC (s-exp->symbol (second l)) (parse (third l)) (parse (fourth l)))]
         [(match) (patternC (parse (second l)) (parse-patterns (rest (rest l))))]))]))
\end{lstlisting}
Three details of the TA solution:
\begin{itemize}
  \item A number clause is \code{cons}ed onto the rest, but a \code{\_} clause returns \rkt|(list (anyP ...))|: parsing \emph{stops} at the first wildcard. Any clause after it could never be reached, so dropping it does not change what the program means: \rkt|(parse '(match 1 (_ 0) (2 3)))| has one clause.
  \item The pattern is \emph{not} parsed as an expression; it must be a number literal or \code{\_}. Anything else is an error, though not one this function raises itself: \code{(match 3 ((+ 1 2) 0))} fails in \code{s-exp->symbol} (\code{not a symbol}), and \code{(match 3 (n 0))} falls off the end of the inner \code{cond} (\code{no matching clause}).
  \item \code{parse-patterns} and \code{parse} call each other, just as the two types refer to each other. \code{match} is a symbol in \emph{our} language, matched by \code{case} like \code{+} and \code{let}.
\end{itemize}

\subsection{Substitution}

\begin{exercise}[Why substitute into \code{match}?]
Why do we need to define substitution for pattern-matching expressions?
\end{exercise}
\begin{solution}
Because \code{let} is implemented by substitution, and the body of a \code{let} can contain a \code{match}. In example~4, the \code{x} and \code{y} bound by the \code{let}s are used in the scrutinees and in the clause bodies. If \code{subst} stopped at a \code{patternC}, those \code{idC}s would never be replaced and \code{eval} would report them as unbound. (In \code{plai-typed} the question does not really arise: \code{type-case} rejects a \code{subst} that has no \code{patternC} case.)
\end{solution}

Substituting into a clause means substituting into its expression; the pattern itself is a number or \code{\_}, with no names in it:
\begin{lstlisting}
(define (subst-pattern (erep : ArithE) (x : symbol)) : (PatternE -> PatternE)
  (lambda (p)
    (type-case PatternE p
      [numP (n e) (numP n (subst erep x e))]
      [anyP (e) (anyP (subst erep x e))])))
\end{lstlisting}
\code{subst-pattern} is \emph{curried}: it takes \code{erep} and \code{x} and returns a function from one clause to one clause. That is exactly the shape \code{map} wants, so the \code{patternC} case is one line:
\begin{lstlisting}
(define (subst (erep : ArithE) (x : symbol) (e : ArithE)) : ArithE
  (type-case ArithE e
    [idC (y) (if (equal? x y) erep (idC y))]
    [numC (n) (numC n)]
    [plusC (e1 e2) (plusC (subst erep x e1) (subst erep x e2))]
    [timesC (e1 e2) (timesC (subst erep x e1) (subst erep x e2))]
    [patternC (e pats) (patternC (subst erep x e) (map (subst-pattern erep x) pats))]
    [letC (y e1 e2)
          (cond
            [(equal? x y) (letC y (subst erep x e1) e2)]
            [else
             (let [(z (gensym))]
               (letC z (subst erep x e1) (subst erep x (subst (idC z) y e2))))])]))
\end{lstlisting}
The \code{letC} case is lecture~4's capture-avoiding \code{subst-lazy-v2}. Note what the \code{patternC} case does \emph{not} need: no shadowing check and no renaming. Those exist for \code{letC} because a \code{let} binds a name; our patterns bind nothing, so a clause body sees exactly the same names as the \code{match} around it, like the children of \code{plusC}.

\subsection{Evaluation}

The helper \code{get-pat-expr} takes the number the scrutinee evaluated to and the clauses, and returns the \emph{expression} of the first clause that matches:
\begin{lstlisting}
(define (get-pat-expr (num : number) (pats : (listof PatternE))) : ArithE
  (cond
    [(empty? pats) (error 'get-pat-expr "unmatched expression")]
    [else (type-case PatternE (first pats)
            [numP (n e) (cond [(equal? num n) e] [else (get-pat-expr num (rest pats))])]
            [anyP (e) e])]))

(define (eval (e : ArithE)) : number
  (type-case ArithE e
    [numC (n) n]
    [plusC (e1 e2) (+ (eval e1) (eval e2))]
    [timesC (e1 e2) (* (eval e1) (eval e2))]
    [patternC (e pats) (eval (get-pat-expr (eval e) pats))]
    [letC (x e1 e2) (eval (subst e1 x e2))]
    [idC (x) (error 'eval "unbound variable!")]))
\end{lstlisting}
\code{eval} evaluates the scrutinee once, picks the clause, and evaluates only that clause's expression.

\subsubsection{Why return the expression, not its value?} \code{get-pat-expr} could evaluate every clause and return numbers. That would be wrong, not just slow: example~4 has a clause \code{(0 z)}, and evaluating it would raise ``unbound variable'' even though the scrutinee is 1 and the clause is never chosen. Returning the unevaluated \tyname{ArithE} means only the chosen clause is evaluated, just as only one branch of an \code{if} is.

\subsubsection{Tracing example 4} \code{g1} stands for whatever name \code{gensym} returns.
\begin{lstlisting}[style=repl]
(let x 1 (let y 1 (match x (0 z) (_ (match y (0 x) (1 (+ 1 y)))))))
letC x:  subst (numC 1) for x in the body.
         The inner let binds y, not x, so y is renamed to g1 first.
---> (let g1 1 (match 1 (0 z) (_ (match g1 (0 1) (1 (+ 1 g1))))))
letC g1: subst (numC 1) for g1 in the body
---> (match 1 (0 z) (_ (match 1 (0 1) (1 (+ 1 1)))))
patternC: scrutinee = 1; (0 z) fails; _ matches
---> (match 1 (0 1) (1 (+ 1 1)))
patternC: scrutinee = 1; (0 1) fails; (1 ...) matches
---> (+ 1 1)
---> 2
\end{lstlisting}
The clause \code{(0 z)} is skipped without being evaluated, which is why the free \code{z} is harmless. Example~3 fails at its first step: the scrutinee is \rkt|(idC 'x)|, no \code{let} has replaced it, and \code{eval} reports it as unbound.

\subsubsection{Clause order} Because \code{get-pat-expr} walks the list from the front and stops at the first match, order matters, and a clause after \code{\_} is dead code:
\begin{lstlisting}
(test (eval (parse '(match 5 (_ 0) (5 1)))) 0)
(test (eval (parse '(match 5 (5 1) (_ 0)))) 1)
(test (eval (parse '(match 5 (5 1) (5 2)))) 1)
(test (eval (parse '(match (* 2 3) (5 0) (6 1) (_ 2)))) 1)
(test/exn (eval (parse '(match 7 (5 0) (6 1)))) "unmatched expression")
\end{lstlisting}
The last test is the run-time ``no matching clause'' error of \Cref{sec:rec3}, now in our own language. Nothing can check for it ahead of time: whether a clause matches depends on a number known only when the program runs.

